\section{Mathematical Modeling Approaches}

This section translates the candidate phenomena into equations you can code.

\subsection{1) Apparent superluminal geometry}

For a pattern moving at \(v=\beta c\):
\begin{equation}
  \beta_{\mathrm{app}} = \frac{\beta\sin\theta}{1-\beta\cos\theta},\qquad v_{\mathrm{app}}=\beta_{\mathrm{app}}c
\end{equation}

Causal constraint: \(\beta<1\). Yet \(\beta_{\mathrm{app}}>1\) is allowed by geometry.

Doppler factor:
\begin{equation}
  \delta = [\Gamma(1-\beta\cos\theta)]^{-1},\qquad \Gamma=(1-\beta^2)^{-1/2}
\end{equation}

Flux scaling (schematic):
\begin{equation}
  I_{\mathrm{obs}}\propto \delta^{p+\alpha}I\,' , \quad p\in\{2,3\}
\end{equation}

\subsection{2) Multi-channel arrival model}

Per source event \(E\), each observable channel has own transport law and delay:
\begin{equation}
  t_i = t_0 + \frac{D}{c} + \tau_{i}(E,\mathbf{M})
\end{equation}

Where \(\tau_i\) includes launch lag, breakout, scattering, diffusion and energy-loss delays.

\subsection{3) Jet power and SMBH coupling (practical form)}

Use a conservative BZ-like parameterization:
\begin{equation}
  P_{\mathrm{BZ}} \approx \kappa\,\frac{\Phi_{\mathrm{BH}}^{2}\Omega_H^2}{4\pi c}
\end{equation}

with horizon spin term \(\Omega_H\propto a_*/M\).

\subsection{4) Shock dynamics}

Non-relativistic self-similar blast:
\begin{equation}
  R=\xi\left(\frac{E}{\rho_0}\right)^{1/5} t^{2/5},\qquad v_s=\frac{2}{5}\frac{R}{t}
\end{equation}

Ultra-relativistic energy-conserving scaling (schematic):
\begin{equation}
  \Gamma^2 R^{3-k}\propto E\quad (\rho\propto r^{-k})
\end{equation}

Transition regime can be handled with moderate-\(\Gamma\) corrections (e.g., Coughlin 2019) for numerical continuity.

\subsection{5) CR transport}

Parker transport form (reduced):
\begin{equation}
  \frac{\partial f}{\partial t}+\mathbf u\cdot\nabla f-\nabla\cdot(\mathbf K\nabla f)-\frac{1}{3}(\nabla\cdot\mathbf u)p\frac{\partial f}{\partial p}=Q
\end{equation}

Simple diffusion scale estimate:
\begin{equation}
  \ell\sim\sqrt{2Kt},\qquad \Delta t\sim \frac{L^2}{2K}
\end{equation}

\subsection{6) Suggested code architecture}

Implement each mechanism as a channel model:
\begin{itemize}
  \item Causal transport module for each channel (GW/EM/CR/MHD).
  \item Geometric projection module for screen-space rendering.
  \item Optional speculative module (off by default).
\end{itemize}

A clean state tuple is:
\begin{equation}
  \mathbf x=\{t,\mathbf r,\Gamma,\theta,\delta, K_{\parallel},K_{\perp}, \tau_{\mathrm{arrive}},\Phi_{\mathrm{BH}},a_*,M_{\mathrm{BH}}\}
\end{equation}

where each field evolves via independent but coordinated update operators.
